\documentclass[../../../include/open-logic-section]{subfiles}

\begin{document}

\olfileid{sfr}{infinite}{card-sb}

\olsection{ضمیمہ: شروڈر-برنسٹین دا ثبوت}

سادی سیٹ تھیوری توں وداع ہون توں پہلاں، ساڈے کول غور کرن لئی اک آخری
سادہ (پر نفیس!) ثبوت اے۔ ایہ شروڈر-برنسٹین
(\olref[sfr][siz][sb]{thm:schroder-bernstein}) دا ثبوت اے: جے
$\cardle{A}{B}$ تے $\cardle{B}{A}$ تاں $\cardeq{A}{B}$؛ یعنی
!!{injection}s $f \colon A \to B$ تے $g \colon B \to A$ دِتے ہوون،
تاں !!a{bijection} $h \colon A \to B$ موجود اے۔

ایس باب وچ اسی ڈیڈیکائنڈ دی \emph{بندشاں} دا تصور اپنایا۔ اصل وچ
ڈیڈیکائنڈ نے ایس تصور نال شروڈر-برنسٹین دا اک سوہنا ثبوت دِتا سی، تے
اسی اوہ ایتھے پیش کراں گے۔ جے تسی ذرا وکھرا پر بنیادی طور تے رلدا ملدا
بیان چاہندے او، تاں ایہ ثبوت \citet[pp.~157--8]{Potter2004} دے بہت
نیڑے چلدا اے۔ گوگل اُتے تھوڑی تلاش وی تہانوں قائل کر دے گی کہ
$\sqrt{2}$ دی غیر ناطقیت وانگ، ایس قضیے دے \emph{کئی} دلچسپ تے وکھرے
ثبوت موجود نیں۔

\olref[sfr][infinite][dedekind]{Closure} ورگی علامت بندی ورتدیاں، حامل
سیٹ C اُتے خود نقشے f تے اوس حامل دے ذیلی سیٹ B لئی رکھو
\[
\Closureofunder{f}{B} = \bigcap \Setabs{X}{B \subseteq X
\text{ تے $X$، $f$-بند اے}}
\]
ہر ایسے سیٹ $B$ تے فنکشن~$f$ لئی۔ ایویں تعریف کیتا
$\Closureofunder{f}{B}$ سب توں چھوٹا $f$-بند سیٹ اے جیہدے وچ~$B$
شامل اے، ایس معنٰی وچ کہ:
% PNBINF-004 ماخذ دی درستی: f نوں حامل C دا خود نقشہ، B نوں C دا ذیلی سیٹ تے X نوں حامل دے ذیلی سیٹاں وچ چلدا سمجھیا اے؛ ماخذ دی بے قید ”ہر فنکشن“ والی عبارت بندش دے وجود تے اگلے اطلاق لئی لازم حامل دائرہ نہیں دیندی۔

\begin{lem}\ollabel{Closureprops}
	ہر ایسے فنکشن $f$ تے ہر $B$ لئی:
	\begin{enumerate}
		\item\ollabel{Closurehaselem} $B \subseteq \Closureofunder{f}{B}$؛ تے
		\item\ollabel{Closureclosed} $\Closureofunder{f}{B}$، $f$-بند اے؛ تے
		\item\ollabel{Closuresmallest} جے $X$، $f$-بند اے تے $B
		\subseteq X$، تاں $\Closureofunder{f}{B} \subseteq X$۔
	\end{enumerate}
\end{lem}

\begin{proof}
عین \olref[sfr][infinite][dedekind]{closureproperties} وانگ۔
\end{proof}

شروڈر-برنسٹین تک پہنچن لئی سانوں اک آخری حقیقت چاہیدی اے:

\begin{prop}\ollabel{sbhelper}
جے $A \subseteq B \subseteq C$ تے $A \approx C$، تاں $\cardeq{\cardeq{A}{B}}{C}$۔
\end{prop}

\begin{proof}
!!a{bijection} $f \colon C \to A$ دِتا ہووے، تاں $F =
\Closureofunder{f}{C \setminus B}$ لو، تے فنکشن $g$، جیہدا دائرۂ
تعریف $C$ اے، دو صورتاں وچ ایویں تعریف کرو:
\[
	g(x) =
	\begin{cases}
			f(x) &\text{جے $x \in F$}\\
			x & \text{نہیں تاں}
	\end{cases}
\]
اسی وکھاواں گے کہ $g$، نقشہ بندی $C \to B$ دے طور تے !!a{bijection} اے؛ ایس توں
$\comp{f^{-1}}{g} \colon A \to B$ دا !!a{bijection} ہونا نکلے گا، تے
ثبوت پورا ہو جاوے گا۔

پہلاں ساڈا دعویٰ اے کہ جے $x \in F$ پر $y\notin F$، تاں $g(x) \neq
g(y)$۔ برہانِ خلف لئی الٹ فرض کرو، سو $y = g(y) = g(x) =
f(x)$۔ کیوں جو $x \in F$ تے \olref{Closureprops} مطابق $F$، $f$-بند
اے، ایس لئی $y = f(x) \in F$، جیہڑا تضاد اے۔

ہُن فرض کرو $g(x) = g(y)$۔ سو اُتلے نتیجے مطابق $x \in F$ اُدوں تے
صرف اُدوں اے جدوں $y \in F$۔ جے $x, y \in F$، تاں $f(x) = g (x) =
g(y) = f(y)$، تے ایس لئی $x = y$، کیوں جو $f$ اک !!a{bijection} اے۔
جے $x, y \notin F$، تاں $x = g(x) = g(y) = y$۔ سو $g$ اک
!!a{injection} اے۔

ہُن صرف ایہ وکھاؤنا رہندا اے کہ $\ran{g} = B$۔ سو $x \in B
\subseteq C$ پکا کرو۔ جے $x \notin F$، تاں $g(x) = x$۔ جے $x \in
F$، تاں $x = f(y)$، کسے $y \in F$ لئی؛ کیوں جو نہیں تاں $F
\setminus \{x\}$، $f$-بند ہوندا تے $C\setminus B$ نوں اپنے وچ شامل رکھدا، جیہڑا
\olref{Closureprops} مطابق ناممکن اے؛ ہُن $g(y) = f(y) = x$۔
\end{proof}

آخر وچ مرکزی نتیجے دا ثبوت ایہ اے۔ فنکشن دے عکس سیٹ لئی یاد کرو کہ
فنکشن $h$ تے سیٹ $D$ دِتے ہوون، تاں اسی $\funimage{h}{D} = \Setabs{h(x)}{x
\in D}$ تعریف کردے آں۔

\begin{proof}[شروڈر-برنسٹین دا ثبوت] $f \colon A \to B$
تے $g \colon B \to A$ نوں !!{injection}s منو۔ کیوں جو $\funimage{f}{A}
\subseteq B$، ایس لئی $\funimage{g}{\funimage{f}{A}} \subseteq
g[B] \subseteq A$۔ نال ای $\comp{f}{g} \colon A \to
\funimage{g}{\funimage{f}{A}}$ اک !!{injection} اے، کیوں جو $g$ تے
$f$ دونیں انجیکشن نیں؛ بلکہ ہدف سیٹ دی ساڈی تعریف دے باعث
$\comp{f}{g}$ اک !!a{bijection} اے۔ سو
$\cardeq{\funimage{g}{\funimage{f}{A}}}{A}$، تے ایس لئی
\olref{sbhelper} مطابق !!a{bijection} $h \colon A \to
\funimage{g}{B}$ موجود اے۔ ہور، $g^{-1}$ اک !!a{bijection}
$\funimage{g}{B} \to B$ اے۔ سو $\comp{h}{g^{-1}} \colon A \to B$
اک !!a{bijection} اے۔
\end{proof}

\end{document}
